  As necessidades dos usuários de internet, atualmente, vem tendo crecimentos que tem utrapassado a capacidade dos provedores de supri-las. Altas variações de requisições em diferentes periodos do dia tem feito com que Provedores de Internet (ISPs) busquem soluções em abordagens que possam facilitar a flexibilidade dos recursos e melhorar o gerenciamento dos mesmos.
  O fatiamento de redes é uma tecnica que visa melhorar a Qualidade de Serviço (Quality of Service - QoS) desses ISPs, particionando recursos e funções em fatias virtuais visando atender essas necessidades tão particulares de cada usuário de rede. Este trabalho apresenta de forma introdutória definições sobre o tema em questão e das abordagens que, em conjunto, tornaram possivel essa evolução. Em seguida apresentamos alguns algoritmos de fatiamento de redes e comparamos seus desempenhos para obter informações mais precisa sobre as diferentes abordagens e em qual cenários cada um desempenha melhor função.